\documentclass[11pt,a4paper]{article}
\usepackage[margin=18mm,top=22mm,bottom=20mm]{geometry}
\usepackage{fontspec}\setmainfont{Latin Modern Roman}
\usepackage{xcolor}\usepackage{amsmath,amssymb,tikz,graphicx,fancyhdr,array,needspace}
\usetikzlibrary{patterns,arrows.meta,decorations.pathmorphing,calc}
\setlength{\headheight}{14pt}
\usetikzlibrary{arrows.meta}
\pagestyle{fancy}\fancyhf{}\renewcommand{\headrulewidth}{0.4pt}
\fancyhead[L]{\textbf{M-02}}\fancyhead[C]{\itshape Fizika — yechimlar (OliyFizika.uz)}\fancyhead[R]{\textbf{Mock №2}}
\fancyfoot[C]{\thepage}
\setlength{\parindent}{0pt}\setlength{\parskip}{3pt}
\newcommand{\opt}[2]{\hspace*{4mm}#1)\ #2\par}
\newcommand{\blank}{\rule{5cm}{0.4pt}}
\begin{document}

\begin{center}{\Large\textbf{FIZIKA — YECHIMLAR}}\\[2pt]{\large Mock №2}\end{center}
\Needspace{10\baselineskip}\noindent\begin{minipage}{\linewidth}
\textbf{1.}\ 3 kg massali jismga 4 m/s$^{2}$ tezlaniş bergan kuç 12 kg massali jismga qanday tezlaniş (m/s$^{2}$) beradi?\par\smallskip
\opt{A}{1}
\opt{B}{3}
\opt{C}{3,5}
\opt{D}{2}
\end{minipage}\vspace{2pt}\par
\noindent\colorbox{gray!12}{\parbox{\dimexpr\linewidth-2\fboxsep}{\textbf{Javob: A) 1}}}\par\smallskip
\noindent\textit{Yechim.}\ Kuç özgarmas: $F=m_1a_1=3\cdot4=12$ N. $a_2=F/m_2=12/12=1$ m/s$^2$.\par\par\medskip\hrule\medskip
\Needspace{10\baselineskip}\noindent\begin{minipage}{\linewidth}
\textbf{2.}\ Qözğalmas blok orqali ötkazilgan ipning bir uçiga 2 kg, ikkinçi uçiga 3 kg massali yuk ilinib tinç holatdan qöyib yuborildi. 3 kg massali jism 2 m pastga tuşganda 2 kg massali jismning kinetik energiyasi (J) qanday böladi? g=10 m/s$^{2}$\par\smallskip
\opt{A}{20}
\opt{B}{12}
\opt{C}{25}
\opt{D}{8}
\end{minipage}\vspace{2pt}\par
\noindent\colorbox{gray!12}{\parbox{\dimexpr\linewidth-2\fboxsep}{\textbf{Javob: D) 8}}}\par\smallskip
\noindent\textit{Yechim.}\ Sistema tezlanişi\par
\[a=\frac{(m_2-m_1)g}{m_1+m_2}=\frac{1\cdot10}{5}=2\ \text{m/s}^2.\]
$s=2$ m dan keyin $v^2=2as=8$ m$^2$/s$^2$, 2 kg li jismning kinetik energiyasi $E_k=\frac{2\cdot8}{2}=8$ J.\par\par\medskip\hrule\medskip
\Needspace{10\baselineskip}\noindent\begin{minipage}{\linewidth}
\textbf{3.}\ Tebranayotgan jismning tebraniş davri 4 s ga teng. Uning çastotasini (Hz) aniqlang.\par\smallskip
\opt{A}{0,25}
\opt{B}{4}
\opt{C}{2}
\opt{D}{0,5}
\end{minipage}\vspace{2pt}\par
\noindent\colorbox{gray!12}{\parbox{\dimexpr\linewidth-2\fboxsep}{\textbf{Javob: A) 0,25}}}\par\smallskip
\noindent\textit{Yechim.}\ $\nu=1/T=1/4=0{,}25$ Hz.\par\par\medskip\hrule\medskip
\Needspace{10\baselineskip}\noindent\begin{minipage}{\linewidth}
\textbf{4.}\ 200 g massali aravaça 1 N kuç ta’sirida gorizontal tekislikda işqalanişsiz tekis tezlanuvçan harakatlanadi. Agar aravaça ustiga 300 g yuk qöysak, huddi şu kuç ta’sirida aravaçaning 2 s da bosib ötadigan yöli neça metrga kamayadi? Ikkala holda ham aravaçaning boşlanğiç tezligi nolga teng.\par\smallskip
\opt{A}{2,5}
\opt{B}{6}
\opt{C}{7}
\opt{D}{8}
\end{minipage}\vspace{2pt}\par
\noindent\colorbox{gray!12}{\parbox{\dimexpr\linewidth-2\fboxsep}{\textbf{Javob: B) 6}}}\par\smallskip
\noindent\textit{Yechim.}\ Yuksiz: $a_1=F/m=1/0{,}2=5$ m/s$^2$, $s_1=a_1t^2/2=10$ m. Yuk bilan: $a_2=1/0{,}5=2$ m/s$^2$, $s_2=4$ m. Yöl $10-4=6$ m ga kamayadi.\par\par\medskip\hrule\medskip
\Needspace{10\baselineskip}\noindent\begin{minipage}{\linewidth}
\textbf{5.}\ 25 m balandlikdan 20 kg massali jism boşlanğiç tezliksiz erkin taşlab yuborildi. Uning kinetik energiyasi 600 J bölganda potensial energiyasi qançaga (J) teng böladi?\par\smallskip
\opt{A}{4400}
\opt{B}{5000}
\opt{C}{3800}
\opt{D}{4200}
\end{minipage}\vspace{2pt}\par
\noindent\colorbox{gray!12}{\parbox{\dimexpr\linewidth-2\fboxsep}{\textbf{Javob: A) 4400}}}\par\smallskip
\noindent\textit{Yechim.}\ Töliq mexanik energiya $E=mgh=20\cdot10\cdot25=5000$ J saqlanadi: $E_p=5000-600=4400$ J.\par\par\medskip\hrule\medskip
\Needspace{10\baselineskip}\noindent\begin{minipage}{\linewidth}
\textbf{6.}\ $p$ va $2p$ impulsli şarlar özaro 60$^\circ$ burçak ostida harakatlanib, noelastik töqnaşdilar va birgalikda harakatlana boşladilar. Şarlarning töqnaşuvdan keyingi impulsi topilsin.\par\smallskip
\opt{A}{$\sqrt{7}p$}
\opt{B}{$2\sqrt{7}p$}
\opt{C}{$3p$}
\opt{D}{$4p$}
\end{minipage}\vspace{2pt}\par
\noindent\colorbox{gray!12}{\parbox{\dimexpr\linewidth-2\fboxsep}{\textbf{Javob: A) $\sqrt{7}p$}}}\par\smallskip
\noindent\textit{Yechim.}\ Impulsning saqlaniş qonuni: $\vec p\,'=\vec p_1+\vec p_2$.\par
\[p'^2=p^2+(2p)^2+2\cdot p\cdot2p\cos60^\circ=7p^2\ \Rightarrow\ p'=\sqrt7\,p.\]\par\medskip\hrule\medskip
\Needspace{10\baselineskip}\noindent\begin{minipage}{\linewidth}
\textbf{7.}\ Jismning mexanik harakatini örganişda nima uçun koordinatalar sistemasidan foydalaniladi?\par\smallskip
\opt{A}{jismning harakatda yoki harakatsiz ekanligini aniqlaş uçun}
\opt{B}{jismning boşlanğiç koordinatalarini aniqlaş uçun}
\opt{C}{jismning fazodagi vaziyatini vaqt ötişi bilan qanday özgarişini aniqlaş uçun}
\opt{D}{jismning boşqa jismlarga nisbatan tinçlikda yoki harakatda ekanligini aniqlaş uçun}
\end{minipage}\vspace{2pt}\par
\noindent\colorbox{gray!12}{\parbox{\dimexpr\linewidth-2\fboxsep}{\textbf{Javob: C) jismning fazodagi vaziyatini vaqt ötişi bilan qanday özgarişini aniqlaş uçun}}}\par\smallskip
\noindent\textit{Yechim.}\ Koordinatalar sistemasi jism vaziyatining vaqt ötişi bilan qanday özgarişini (harakat qonunini) aniqlaş uçun kerak.\par\par\medskip\hrule\medskip
\Needspace{10\baselineskip}\noindent\begin{minipage}{\linewidth}
\textbf{8.}\ Massasi M bölgan silindrga ip öralgan. Söngra silindr pastga taşlab yuborilib, ip yuqoriga tortib turiladi. Bunda silindrning massa markazi ipning yoyilişi davrida aynan bir joyda qoldi. Ipning taranglik kuçini aniqlang. g-erkin tuşiş tezlanişi.\par\smallskip
\opt{A}{Mg}
\opt{B}{Mg/2}
\opt{C}{2Mg}
\opt{D}{0}
\end{minipage}\vspace{2pt}\par
\noindent\colorbox{gray!12}{\parbox{\dimexpr\linewidth-2\fboxsep}{\textbf{Javob: A) Mg}}}\par\smallskip
\noindent\textit{Yechim.}\ Massa markazi joyida qolsa, uning tezlanişi nol: $T-Mg=0$, ya’ni $T=Mg$.\par\par\medskip\hrule\medskip
\Needspace{10\baselineskip}\noindent\begin{minipage}{\linewidth}
\textbf{9.}\ Aylanma harakat qilayotgan jism 3 minutda 3$\pi$ radianga burildi. Uning tebraniş davrini (s) aniqlang.\par\smallskip
\opt{A}{100}
\opt{B}{60}
\opt{C}{120}
\opt{D}{180}
\end{minipage}\vspace{2pt}\par
\noindent\colorbox{gray!12}{\parbox{\dimexpr\linewidth-2\fboxsep}{\textbf{Javob: C) 120}}}\par\smallskip
\noindent\textit{Yechim.}\ Burçak tezligi $\omega=\varphi/t=3\pi/180=\pi/60$ rad/s.\par
\[T=\frac{2\pi}{\omega}=\frac{2\pi}{\pi/60}=120\ \text{s}.\]\par\medskip\hrule\medskip
\Needspace{10\baselineskip}\noindent\begin{minipage}{\linewidth}
\textbf{10.}\ Öq uzgan mergan 2,7 s ötgaç öqning nişonga tekkanida çiqqan tovuşni eşitdi. Öqning tezligi 800 m/s, tovuşning havodagi tezligi 330 m/s bölsa, nişongaça bölgan masofa (m) topilsin.\par\smallskip
\opt{A}{690}
\opt{B}{775}
\opt{C}{630}
\opt{D}{720}
\end{minipage}\vspace{2pt}\par
\noindent\colorbox{gray!12}{\parbox{\dimexpr\linewidth-2\fboxsep}{\textbf{Javob: C) 630}}}\par\smallskip
\noindent\textit{Yechim.}\ Öqning uçiş vaqti va tovuşning qaytiş vaqti yiğindisi: $\frac{s}{800}+\frac{s}{330}=2{,}7$.\par
\[s=\frac{2{,}7\cdot800\cdot330}{1130}\approx630\ \text{m}.\]\par\medskip\hrule\medskip
\Needspace{10\baselineskip}\noindent\begin{minipage}{\linewidth}
\textbf{11.}\ 1 mol miqdordagi vodorod va geliy gazlarining bir xil temperaturadagi içki energiyalarini soliştiring.\par\smallskip
\opt{A}{geliyniki katta}
\opt{B}{aniqlab bölmaydi}
\opt{C}{vodorodniki katta}
\opt{D}{bir xil}
\end{minipage}\vspace{2pt}\par
\noindent\colorbox{gray!12}{\parbox{\dimexpr\linewidth-2\fboxsep}{\textbf{Javob: C) vodorodniki katta}}}\par\smallskip
\noindent\textit{Yechim.}\ $U=\frac{i}{2}\nu RT$. Vodorod ikki atomli ($i=5$), geliy bir atomli ($i=3$), şuning uçun vodorodning içki energiyasi katta.\par\par\medskip\hrule\medskip
\Needspace{10\baselineskip}\noindent\begin{minipage}{\linewidth}
\textbf{12.}\ Kapillar simobga tuşurilgan, simobning sirt taranglik koeffitsiyenti 0,51 N/m, ziçligi $13600\ \mathrm{kg/m^3}$. Simob kapillarni höllamagani uçun kapillar içidagi simob sathi atrofdagi sathdan 5 mm pastroq. Kapillarning içki radiusi (mm) qanday? $g=10\ \mathrm{m/s^2}$\par\smallskip
\opt{A}{0,19}
\opt{B}{1,5}
\opt{C}{0,37}
\opt{D}{0,75}
\end{minipage}\vspace{2pt}\par
\noindent\colorbox{gray!12}{\parbox{\dimexpr\linewidth-2\fboxsep}{\textbf{Javob: B) 1,5}}}\par\smallskip
\noindent\textit{Yechim.}\ Kapillyardagi sath özgarişi $h=\frac{2\sigma}{\rho gr}$:\par
\[r=\frac{2\sigma}{\rho gh}=\frac{2\cdot0{,}51}{13600\cdot10\cdot0{,}005}=1{,}5\cdot10^{-3}\ \text{m}=1{,}5\ \text{mm}.\]\par\medskip\hrule\medskip
\Needspace{10\baselineskip}\noindent\begin{minipage}{\linewidth}
\textbf{13.}\ Gaz adiabatik kengayib 24 kJ iş bajargan bölsa, uning içki energiyasi qanday özgaradi?\par\smallskip
\opt{A}{24 kJ ga kamayadi}
\opt{B}{24 kJ ga ortadi}
\opt{C}{özgarmaydi}
\opt{D}{12 kJ ga kamayadi}
\end{minipage}\vspace{2pt}\par
\noindent\colorbox{gray!12}{\parbox{\dimexpr\linewidth-2\fboxsep}{\textbf{Javob: A) 24 kJ ga kamayadi}}}\par\smallskip
\noindent\textit{Yechim.}\ Adiabatik jarayonda $Q=0$: $\Delta U=-A=-24$ kJ, ya’ni içki energiya 24 kJ ga kamayadi.\par\par\medskip\hrule\medskip
\Needspace{10\baselineskip}\noindent\begin{minipage}{\linewidth}
\textbf{14.}\ $27\ ^\circ$C haroratli vodorod molekulalarining örtaça kvadratik tezligini (m/s) va şu tezlikdagi bitta molekulaning impulsini toping (kg$\cdot$m/s). Vodorodning molyar massasi 2 g/mol.\par\smallskip
\opt{A}{$1933;\ 6{,}4\cdot10^{-24}$}
\opt{B}{$1733;\ 5{,}64\cdot10^{-24}$}
\opt{C}{$1033;\ 6{,}44\cdot10^{-24}$}
\opt{D}{$193;\ 6{,}64\cdot10^{-24}$}
\end{minipage}\vspace{2pt}\par
\noindent\colorbox{gray!12}{\parbox{\dimexpr\linewidth-2\fboxsep}{\textbf{Javob: A) $1933;\ 6{,}4\cdot10^{-24}$}}}\par\smallskip
\noindent\textit{Yechim.}\ \[v=\sqrt{\frac{3RT}{M}}=\sqrt{\frac{3\cdot8{,}31\cdot300}{0{,}002}}\approx1933\ \text{m/s}.\]
Bitta molekula massasi $m_0=M/N_A=0{,}002/(6{,}02\cdot10^{23})\approx3{,}32\cdot10^{-27}$ kg, impulsi $p=m_0v\approx6{,}4\cdot10^{-24}$ kg$\cdot$m/s.\par\par\medskip\hrule\medskip
\Needspace{10\baselineskip}\noindent\begin{minipage}{\linewidth}
\textbf{15.}\ Gaz adiabatik siqilganda temperatura 40 K ga ortdi. Agar bu gaz 320 g kislorod ekanligi ma’lum bölsa, bunda qança iş bajarilgan (J)?\par\smallskip
\opt{A}{4986}
\opt{B}{3324}
\opt{C}{1243}
\opt{D}{8310}
\end{minipage}\vspace{2pt}\par
\noindent\colorbox{gray!12}{\parbox{\dimexpr\linewidth-2\fboxsep}{\textbf{Javob: D) 8310}}}\par\smallskip
\noindent\textit{Yechim.}\ Adiabatik siqişda taşqi kuçlar işi içki energiya özgarişiga teng. Kislorod ikki atomli, $\nu=320/32=10$ mol:\par
\[A=\Delta U=\frac52\nu R\Delta T=\frac52\cdot10\cdot8{,}31\cdot40=8310\ \text{J}.\]\par\medskip\hrule\medskip
\Needspace{10\baselineskip}\noindent\begin{minipage}{\linewidth}
\textbf{16.}\ Qirrasi $a$ ga teng bölgan tetraedr şaklidagi jism sirti böylab $q$ zaryad bir tekisda taqsimlangan. Zaryadning sirtiy ziçligini aniqlang.\par\smallskip
\opt{A}{$\dfrac{q}{a^2}$}
\opt{B}{$\dfrac{q}{a\sqrt3}$}
\opt{C}{$\dfrac{q}{a^3}$}
\opt{D}{$\dfrac{q}{a^2\sqrt3}$}
\end{minipage}\vspace{2pt}\par
\noindent\colorbox{gray!12}{\parbox{\dimexpr\linewidth-2\fboxsep}{\textbf{Javob: D) $\dfrac{q}{a^2\sqrt3}$}}}\par\smallskip
\noindent\textit{Yechim.}\ Tetraedrning 4 ta yoği — tomoni $a$ bölgan teng tomonli uçburçaklar: $S=4\cdot\frac{\sqrt3}{4}a^2=\sqrt3a^2$. $\sigma=\frac{q}{S}=\frac{q}{a^2\sqrt3}$.\par\par\medskip\hrule\medskip
\Needspace{10\baselineskip}\noindent\begin{minipage}{\linewidth}
\textbf{17.}\ Massalari $m$ dan bölgan $q$ va $3q$ zaryadli zarralar bir biridan $a$ masofada uşlab turilibdi. $q$ zaryadli zarra böşatib yuborilsa, u qanday maksimal tezlikka erişadi?\par\smallskip
\opt{A}{$\sqrt{\dfrac{3q^2}{am}}$}
\opt{B}{$\sqrt{\dfrac{6kq^2}{am}}$}
\opt{C}{$\sqrt{\dfrac{12kq}{am}}$}
\opt{D}{$\sqrt{\dfrac{18kq}{am}}$}
\end{minipage}\vspace{2pt}\par
\noindent\colorbox{gray!12}{\parbox{\dimexpr\linewidth-2\fboxsep}{\textbf{Javob: B) $\sqrt{\dfrac{6kq^2}{am}}$}}}\par\smallskip
\noindent\textit{Yechim.}\ $3q$ zaryad qözğalmas, $q$ zaryadli zarra çeksiz uzoqlaşadi; özaro ta’sir energiyasi kinetik energiyaga aylanadi:\par
\[\frac{k\cdot q\cdot3q}{a}=\frac{mv^2}{2}\ \Rightarrow\ v=\sqrt{\frac{6kq^2}{am}}.\]\par\medskip\hrule\medskip
\Needspace{10\baselineskip}\noindent\begin{minipage}{\linewidth}
\textbf{18.}\ 50 Hz çastotali özgaruvçan tok manbayiga induktivligi 25 mH bölgan ğaltak ulangan. Agar zanjirdagi tokning maksimal qiymati 2 A bölsa, zanjirdagi kuçlanişning maksimal qiymatini (V) aniqlang.\par\smallskip
\opt{A}{2,5}
\opt{B}{5}
\opt{C}{5$\pi$}
\opt{D}{2,5$\pi$}
\end{minipage}\vspace{2pt}\par
\noindent\colorbox{gray!12}{\parbox{\dimexpr\linewidth-2\fboxsep}{\textbf{Javob: C) 5$\pi$}}}\par\smallskip
\noindent\textit{Yechim.}\ $X_L=\omega L=2\pi\nu L=2\pi\cdot50\cdot0{,}025=2{,}5\pi\ \Omega$. $U_0=I_0X_L=2\cdot2{,}5\pi=5\pi$ V.\par\par\medskip\hrule\medskip
\Needspace{10\baselineskip}\noindent\begin{minipage}{\linewidth}
\textbf{19.}\ Radiusi 1 mm bölgan silindrik mis sim orqali 3,14 A tok ötmoqda. Tok ziçligini aniqlang (A/m$^{2}$).\par\smallskip
\opt{A}{$2\cdot10^5$}
\opt{B}{$6{,}28\cdot10^5$}
\opt{C}{$1\cdot10^6$}
\opt{D}{$3{,}14\cdot10^6$}
\end{minipage}\vspace{2pt}\par
\noindent\colorbox{gray!12}{\parbox{\dimexpr\linewidth-2\fboxsep}{\textbf{Javob: C) $1\cdot10^6$}}}\par\smallskip
\noindent\textit{Yechim.}\ \[j=\frac{I}{\pi r^2}=\frac{3{,}14}{3{,}14\cdot10^{-6}}=10^6\ \text{A/m}^2.\]\par\medskip\hrule\medskip
\Needspace{10\baselineskip}\noindent\begin{minipage}{\linewidth}
\textbf{20.}\ Bir xil q zaryadga ega bölgan 2 cm va 4 cm radiusli şarlar bir-biriga tekkizilib, yana avvalgi masofaga uzoqlaştirildi. Bunda şarlarning özaro ta’sirlaşiş kuçi qanday özgaradi?\par\smallskip
\opt{A}{9/8 marta kamayadi}
\opt{B}{9/8 marta ortadi}
\opt{C}{4/3 marta ortadi}
\opt{D}{özgarmaydi}
\end{minipage}\vspace{2pt}\par
\noindent\colorbox{gray!12}{\parbox{\dimexpr\linewidth-2\fboxsep}{\textbf{Javob: A) 9/8 marta kamayadi}}}\par\smallskip
\noindent\textit{Yechim.}\ Tekkizilganda umumiy zaryad $2q$ radiuslarga proporsional taqsimlanadi (potensiallar tenglaşadi): $q_1=\frac{2q\cdot2}{6}=\frac{2q}{3}$, $q_2=\frac{4q}{3}$. Kuç $q_1q_2=\frac{8}{9}q^2$ ga proporsional, ya’ni $9/8$ marta kamayadi.\par\par\medskip\hrule\medskip
\Needspace{10\baselineskip}\noindent\begin{minipage}{\linewidth}
\textbf{21.}\ Özgarmas kuçlaniş manbaiga ulangan kondensator qoplamalari orasidagi masofa 3 marta orttirilsa, elektr maydon energiya ziçligi qanday özgaradi?\par\smallskip
\opt{A}{9 marta kamayadi}
\opt{B}{9 marta ortadi}
\opt{C}{3 marta kamayadi}
\opt{D}{özgarmaydi}
\end{minipage}\vspace{2pt}\par
\noindent\colorbox{gray!12}{\parbox{\dimexpr\linewidth-2\fboxsep}{\textbf{Javob: A) 9 marta kamayadi}}}\par\smallskip
\noindent\textit{Yechim.}\ $U$ özgarmas, $E=U/d$ 3 marta kamayadi, $w=\varepsilon_0E^2/2$ esa 9 marta kamayadi.\par\par\medskip\hrule\medskip
\Needspace{10\baselineskip}\noindent\begin{minipage}{\linewidth}
\textbf{22.}\ Maksimal tok kuçi 5 A bölgan özgaruvçan tok zanjiridagi tokning effektiv qiymatini (A) aniqlang.\par\smallskip
\opt{A}{$\dfrac{5}{\sqrt{2}}$}
\opt{B}{$\dfrac{3}{\sqrt{2}}$}
\opt{C}{5}
\opt{D}{$\dfrac{5}{\sqrt{3}}$}
\end{minipage}\vspace{2pt}\par
\noindent\colorbox{gray!12}{\parbox{\dimexpr\linewidth-2\fboxsep}{\textbf{Javob: A) $\dfrac{5}{\sqrt{2}}$}}}\par\smallskip
\noindent\textit{Yechim.}\ $I_{ef}=I_0/\sqrt2=5/\sqrt2$ A.\par\par\medskip\hrule\medskip
\Needspace{10\baselineskip}\noindent\begin{minipage}{\linewidth}
\textbf{23.}\ Elektr maydonda 1 $\mu$g massali çang zarrasi vertikal ravişda 4 sekundda yerga tuşdi. Agar maydon kuçlanganligi 50 V/m va maydon kuç çiziqlari tik yuqoriga yönalgan bölsa, çang zarrasi qanday balandlikdan tuşgan (m)? Çang zarrasi zaryadi 0,02 nC.\par\smallskip
\opt{A}{80}
\opt{B}{63}
\opt{C}{79}
\opt{D}{72}
\end{minipage}\vspace{2pt}\par
\noindent\colorbox{gray!12}{\parbox{\dimexpr\linewidth-2\fboxsep}{\textbf{Javob: D) 72}}}\par\smallskip
\noindent\textit{Yechim.}\ $mg=10^{-9}\cdot10=10^{-8}$ N, elektr kuçi $qE=2\cdot10^{-11}\cdot50=10^{-9}$ N yuqoriga.\par
\[a=\frac{mg-qE}{m}=\frac{10^{-8}-10^{-9}}{10^{-9}}=9\ \text{m/s}^2,\qquad h=\frac{at^2}{2}=\frac{9\cdot16}{2}=72\ \text{m}.\]\par\medskip\hrule\medskip
\Needspace{10\baselineskip}\noindent\begin{minipage}{\linewidth}
\textbf{24.}\ Har birining zaryadi $q$ ga teng, radiuslari esa $r$ va $R$ bölgan ikki ötkazgiç şarlar bir-biriga tekkizilib uzoqlaştirilsa, birinçisining zaryadi qanday bölib qoladi?\par\smallskip
\opt{A}{$\dfrac{qr}{r+R}$}
\opt{B}{$\dfrac{2qR}{r+R}$}
\opt{C}{$\dfrac{qR}{r+R}$}
\opt{D}{$\dfrac{2qr}{r+R}$}
\end{minipage}\vspace{2pt}\par
\noindent\colorbox{gray!12}{\parbox{\dimexpr\linewidth-2\fboxsep}{\textbf{Javob: D) $\dfrac{2qr}{r+R}$}}}\par\smallskip
\noindent\textit{Yechim.}\ Tekkizilganda potensiallar tenglaşadi: $q_1/r=q_2/R$, $q_1+q_2=2q$. Bundan $q_1=\frac{2qr}{r+R}$.\par\par\medskip\hrule\medskip
\Needspace{10\baselineskip}\noindent\begin{minipage}{\linewidth}
\textbf{25.}\ Potensiallar farqi 4 V bölgan elektr maydonda 8 mC zaryadli zarraçani köçirişda elektr maydon qança iş (mJ) bajaradi?\par\smallskip
\opt{A}{32}
\opt{B}{16}
\opt{C}{64}
\opt{D}{2}
\end{minipage}\vspace{2pt}\par
\noindent\colorbox{gray!12}{\parbox{\dimexpr\linewidth-2\fboxsep}{\textbf{Javob: A) 32}}}\par\smallskip
\noindent\textit{Yechim.}\ $A=qU=8\cdot10^{-3}\cdot4=32\cdot10^{-3}$ J $=32$ mJ.\par\par\medskip\hrule\medskip
\Needspace{10\baselineskip}\noindent\begin{minipage}{\linewidth}
\textbf{26.}\ Botiq közguning egrilik radiusi 24 cm. Uning fokus masofasini (cm) aniqlang.\par\smallskip
\opt{A}{24}
\opt{B}{6}
\opt{C}{12}
\opt{D}{48}
\end{minipage}\vspace{2pt}\par
\noindent\colorbox{gray!12}{\parbox{\dimexpr\linewidth-2\fboxsep}{\textbf{Javob: C) 12}}}\par\smallskip
\noindent\textit{Yechim.}\ $F=R/2=24/2=12$ cm.\par\par\medskip\hrule\medskip
\Needspace{10\baselineskip}\noindent\begin{minipage}{\linewidth}
\textbf{27.}\ Yassi közgu va odam orasidagi masofa 4 m ga teng. Odam közgudan 1 m ga uzoqlaşdi. Közgu ham xuddi şu yönalişda 2 m ga surildi. Odam va uning tasviri orasidagi masofa qanday bölib qoladi (m)?\par\smallskip
\opt{A}{8}
\opt{B}{3}
\opt{C}{4}
\opt{D}{6}
\end{minipage}\vspace{2pt}\par
\noindent\colorbox{gray!12}{\parbox{\dimexpr\linewidth-2\fboxsep}{\textbf{Javob: D) 6}}}\par\smallskip
\noindent\textit{Yechim.}\ Odam közgudan $4+1=5$ m uzoqlaşdi, közgu esa odam tomon 2 m surildi: oraliq $5-2=3$ m. Tasvir közgu ortida xuddi şunday masofada, demak odam va tasvir orasidagi masofa $2\cdot3=6$ m.\par\par\medskip\hrule\medskip
\Needspace{10\baselineskip}\noindent\begin{minipage}{\linewidth}
\textbf{28.}\ Optik kuçi 4 dptr bölgan yupqa linzaning fokus masofasini (cm) aniqlang.\par\smallskip
\opt{A}{25}
\opt{B}{20}
\opt{C}{2,2}
\opt{D}{2}
\end{minipage}\vspace{2pt}\par
\noindent\colorbox{gray!12}{\parbox{\dimexpr\linewidth-2\fboxsep}{\textbf{Javob: A) 25}}}\par\smallskip
\noindent\textit{Yechim.}\ $F=1/D=1/4=0{,}25$ m $=25$ cm.\par\par\medskip\hrule\medskip
\Needspace{10\baselineskip}\noindent\begin{minipage}{\linewidth}
\textbf{29.}\ Yadro reaktoridagi katta energiyalik $^7_3Li$ ning yadrosi nömalum $X$ zarra yadrosi bilan reaksiyaga kirişişi natijasida ikkita alfa-zarra hosil böldi. $X$ zarrani aniqlang.\par\smallskip
\opt{A}{$^3_1H$}
\opt{B}{$^1_1H$}
\opt{C}{$^4_2He$}
\opt{D}{$^2_1H$}
\end{minipage}\vspace{2pt}\par
\noindent\colorbox{gray!12}{\parbox{\dimexpr\linewidth-2\fboxsep}{\textbf{Javob: B) $^1_1H$}}}\par\smallskip
\noindent\textit{Yechim.}\ $^7_3\mathrm{Li}+^A_Z\mathrm X\to2\,^4_2\mathrm{He}$. Zaryad: $3+Z=4$, $Z=1$; massa soni: $7+A=8$, $A=1$. Demak $X$ — proton $^1_1\mathrm H$.\par\par\medskip\hrule\medskip
\Needspace{10\baselineskip}\noindent\begin{minipage}{\linewidth}
\textbf{30.}\ Tezligi 0,1$c$ bölgan zarraning impulsi $p$ ga teng. Tezligi 2 marta ortsa, zarraning töliq energiyasi qanday böladi? $c$ - yoruğlik tezligi.\par\smallskip
\opt{A}{$\sqrt{119}pc$}
\opt{B}{$\sqrt{132}pc$}
\opt{C}{$\sqrt{90}pc$}
\opt{D}{$\sqrt{103}pc$}
\end{minipage}\vspace{2pt}\par
\noindent\colorbox{gray!12}{\parbox{\dimexpr\linewidth-2\fboxsep}{\textbf{Javob: D) $\sqrt{103}pc$}}}\par\smallskip
\noindent\textit{Yechim.}\ $p=\frac{mv}{\sqrt{1-v^2/c^2}}$: $v=0{,}1c$ da $p=\frac{0{,}1mc}{\sqrt{0{,}99}}$, bundan $mc^2=\frac{10\sqrt{0{,}99}}{1}\,pc$. $v=0{,}2c$ da\par
\[E=\frac{mc^2}{\sqrt{0{,}96}}=10\sqrt{\frac{0{,}99}{0{,}96}}\,pc=\sqrt{103{,}1}\,pc\approx\sqrt{103}\,pc.\]\par\medskip\hrule\medskip
\Needspace{10\baselineskip}\noindent\begin{minipage}{\linewidth}
\textbf{31.}\ Radioaktiv preparatning yarim yemiriliş davri 4 minut, radioaktiv atomlari soni $10^7$ ta. 12 minut içida qança radiaktiv atom parçalanadi?\par\smallskip
\opt{A}{$6{,}5\cdot10^6$}
\opt{B}{$9{,}1\cdot10^6$}
\opt{C}{$9{,}25\cdot10^6$}
\opt{D}{$8{,}75\cdot10^6$}
\end{minipage}\vspace{2pt}\par
\noindent\colorbox{gray!12}{\parbox{\dimexpr\linewidth-2\fboxsep}{\textbf{Javob: D) $8{,}75\cdot10^6$}}}\par\smallskip
\noindent\textit{Yechim.}\ 12 minut — 3 ta yarim yemiriliş davri: $N=N_0/8$. Parçalangan atomlar\par
\[\Delta N=N_0\left(1-\frac18\right)=\frac78\cdot10^7=8{,}75\cdot10^6.\]\par\medskip\hrule\medskip
\Needspace{10\baselineskip}\noindent\begin{minipage}{\linewidth}
\textbf{32.}\ Jism sinus qonuniga böysungan holda tebranmoqda. Tebraniş boşidan hisoblaganda qanday eng qisqa vaqtdan keyin jism impulsi 2 marta kamayadi? Jismning tebraniş davri T ga teng.\par\smallskip
\opt{A}{T/12}
\opt{B}{T/8}
\opt{C}{T/3}
\opt{D}{T/6}
\end{minipage}\vspace{2pt}\par
\noindent\colorbox{gray!12}{\parbox{\dimexpr\linewidth-2\fboxsep}{\textbf{Javob: D) T/6}}}\par\smallskip
\noindent\textit{Yechim.}\ $x=A\sin\omega t$ bölsa, $v=A\omega\cos\omega t$; boşlanişda impuls maksimal. Impuls 2 marta kamayişi uçun $\cos\omega t=\frac12$, $\omega t=\frac{\pi}{3}$:\par
\[t=\frac{\pi/3}{2\pi/T}=\frac{T}{6}.\]\par\medskip\hrule\medskip
\newpage\noindent\textbf{33–35.}\ Rasmda elektr zanjir tasvirlangan: $\varepsilon=40$ V, $r=1\ \Omega$, $R_1=8\ \Omega$, $R_2=8\ \Omega$, $R_3=5\ \Omega$.\par\medskip
\begin{center}\begin{tikzpicture}[x=1cm,y=1cm,line width=0.5pt,>=Stealth,font=\small]
\draw (0,0.7)--(0,0)--(5.2,0)--(5.2,2.0)--(4.6,2.0); \draw (0,1.0)--(0,2.0)--(1.2,2.0);
\draw (-0.3,1.0)--(0.3,1.0); \draw (-0.15,0.7)--(0.15,0.7); \node[left] at (-0.35,0.85) {$\varepsilon,\ r$};
\draw (1.2,2.9)--(1.2,1.1) (3.2,2.9)--(3.2,1.1);
\draw (1.2,2.9)--(1.8,2.9); \draw[fill=white] (1.8,2.75) rectangle (2.6,3.05); \node[above] at (2.2,3.05) {$R_1$}; \draw (2.6,2.9)--(3.2,2.9);
\draw (1.2,1.1)--(1.8,1.1); \draw[fill=white] (1.8,0.95) rectangle (2.6,1.25); \node[above] at (2.2,1.25) {$R_2$}; \draw (2.6,1.1)--(3.2,1.1);
\draw (3.2,2.0)--(3.8,2.0); \draw[fill=white] (3.8,1.85) rectangle (4.6,2.15); \node[above] at (4.2,2.15) {$R_3$};
\fill (1.2,2.0) circle(1.2pt) (3.2,2.0) circle(1.2pt);
\end{tikzpicture}\end{center}
\noindent\textbf{33.}\ Batareyadan ötayotgan tok kuçini (A) aniqlang.\par
\noindent\textbf{34.}\ Batareya qisqiçlaridagi kuçlanişni (V) toping.\par
\noindent\textbf{35.}\ $R_1$ qarşilikdan ötayotgan tok kuçini (A) aniqlang.\par
\noindent Variantlar: A) 2;\quad B) 4;\quad C) 5;\quad D) 8;\quad E) 32;\quad F) 36\par\smallskip
\noindent\colorbox{gray!12}{\parbox{\dimexpr\linewidth-2\fboxsep}{\textbf{Javob: 33--B) 4; 34--F) 36; 35--A) 2}}}\par\smallskip
\noindent\textit{Yechim.}\ $R_1\parallel R_2$: $R_{12}=4\ \Omega$. Taşqi qarşilik $R=4+5=9\ \Omega$.\par
\[I=\frac{\varepsilon}{R+r}=\frac{40}{10}=4\ \text{A},\qquad U=\varepsilon-Ir=40-4=36\ \text{V}.\]
$R_1=R_2$ bölgani uçun tok teng bölinadi: $I_1=2$ A.\par\par\medskip\hrule\medskip
\Needspace{12\baselineskip}\noindent\textbf{36.}\ Ziçligi 600 kg/m$^3$ bölgan, qirrasi 10 cm li yoğoç kub suvda suzib yuribdi. $\rho_{suv}=1000$ kg/m$^3$, $g=10$ m/s$^2$.\par
\noindent\hspace*{4mm}a)\ Kubning suvga botgan qismi çuqurligini (cm) toping.\par
\noindent\hspace*{4mm}b)\ Kubni suvga töliq botiriş uçun unga qöyiş kerak bölgan eng kiçik vertikal kuçni (N) toping.\par
\noindent\colorbox{gray!12}{\parbox{\dimexpr\linewidth-2\fboxsep}{\textbf{Javob: a)\ 6;\quad b)\ 4}}}\par\smallskip
\noindent\textit{Yechim.}\ Suziş şarti: $\rho_s ghS=\rho gaS$, botgan çuqurlik $h=\frac{\rho}{\rho_s}a=0{,}6\cdot10=6$ cm. Töliq botiriş uçun qöşimça kuç $F=(\rho_s-\rho)gV=400\cdot10\cdot10^{-3}=4$ N.\par\par\medskip\hrule\medskip
\Needspace{12\baselineskip}\noindent\textbf{37.}\ Jism 80 m balandlikdan boşlanğiç tezliksiz erkin tuşa boşladi. $g=10$ m/s$^2$.\par
\noindent\hspace*{4mm}a)\ Jismning tuşiş vaqtini (s) toping.\par
\noindent\hspace*{4mm}b)\ Jismning harakatining oxirgi sekundida bosib ötgan yölini (m) toping.\par
\noindent\colorbox{gray!12}{\parbox{\dimexpr\linewidth-2\fboxsep}{\textbf{Javob: a)\ 4;\quad b)\ 35}}}\par\smallskip
\noindent\textit{Yechim.}\ $t=\sqrt{2h/g}=\sqrt{16}=4$ s. Dastlabki 3 s da $gt^2/2=45$ m ötadi, oxirgi sekundda $80-45=35$ m.\par\par\medskip\hrule\medskip
\Needspace{12\baselineskip}\noindent\textbf{38.}\ Massasi 10 g bölgan öq 400 m/s tezlik bilan uçib kelib, ipga osilgan 1,99 kg massali qum töldirilgan qutiga tiqilib qoldi. $g=10$ m/s$^2$.\par
\noindent\hspace*{4mm}a)\ Öq tiqilgandan keyin qutining tezligini (m/s) toping.\par
\noindent\hspace*{4mm}b)\ Quti qanday balandlikka (cm) kötariladi?\par
\noindent\colorbox{gray!12}{\parbox{\dimexpr\linewidth-2\fboxsep}{\textbf{Javob: a)\ 2;\quad b)\ 20}}}\par\smallskip
\noindent\textit{Yechim.}\ Impulsning saqlaniş qonuni (noelastik uriliş):\par
\[u=\frac{mv}{m+M}=\frac{0{,}01\cdot400}{2}=2\ \text{m/s},\qquad h=\frac{u^2}{2g}=\frac{4}{20}=0{,}2\ \text{m}=20\ \text{cm}.\]\par\medskip\hrule\medskip
\Needspace{12\baselineskip}\noindent\textbf{39.}\ Hajmi 8,31 L bölgan ballonda 27 $^\circ$C temperaturada 64 g kislorod bor. $M=32$ g/mol, $R=8{,}31$ J/(mol$\cdot$K).\par
\noindent\hspace*{4mm}a)\ Ballondagi gaz bosimini (kPa) toping.\par
\noindent\hspace*{4mm}b)\ Ballon 127 $^\circ$C gaça qizdirilganda gaz bosimi qanday böladi (kPa)?\par
\noindent\colorbox{gray!12}{\parbox{\dimexpr\linewidth-2\fboxsep}{\textbf{Javob: a)\ 600;\quad b)\ 800}}}\par\smallskip
\noindent\textit{Yechim.}\ $\nu=64/32=2$ mol.\par
\[p_1=\frac{\nu RT}{V}=\frac{2\cdot8{,}31\cdot300}{8{,}31\cdot10^{-3}}=6\cdot10^5\ \text{Pa}=600\ \text{kPa}.\]
Izoxorik jarayon: $p_2=p_1T_2/T_1=600\cdot400/300=800$ kPa.\par\par\medskip\hrule\medskip
\Needspace{12\baselineskip}\noindent\textbf{40.}\ 80 $^\circ$C temperaturali 2 kg suvga 20 $^\circ$C temperaturali 3 kg suv quyildi. Issiqlik isroflari hisobga olinmasin. $c_{suv}=4200$ J/(kg$\cdot$K).\par
\noindent\hspace*{4mm}a)\ Aralaşmaning qaror topgan temperaturasini ($^\circ$C) toping.\par
\noindent\hspace*{4mm}b)\ Issiq suv bergan issiqlik miqdorini (kJ) toping.\par
\noindent\colorbox{gray!12}{\parbox{\dimexpr\linewidth-2\fboxsep}{\textbf{Javob: a)\ 44;\quad b)\ 302,4}}}\par\smallskip
\noindent\textit{Yechim.}\ Issiqlik balansi: $2(80-t)=3(t-20)$, $t=44\ ^\circ$C. Issiq suv bergan issiqlik $Q=4200\cdot2\cdot36=302\,400$ J $=302{,}4$ kJ.\par\par\medskip\hrule\medskip
\Needspace{12\baselineskip}\noindent\textbf{41.}\ Uzunligi 0,5 m bölgan töğri ötkazgiçdan 4 A tok ötmoqda. Ötkazgiç induksiyasi 0,2 T bölgan bir jinsli magnit maydonda induksiya çiziqlariga $30^\circ$ burçak ostida joylaşgan.\par
\noindent\hspace*{4mm}a)\ Ötkazgiçga ta’sir qiluvçi Amper kuçini (N) toping.\par
\noindent\hspace*{4mm}b)\ Ötkazgiçni qanday burab qöyganda kuç eng katta böladi, şu eng katta kuçni (N) toping.\par
\noindent\colorbox{gray!12}{\parbox{\dimexpr\linewidth-2\fboxsep}{\textbf{Javob: a)\ 0,2;\quad b)\ 0,4}}}\par\smallskip
\noindent\textit{Yechim.}\ $F=BIL\sin\alpha=0{,}2\cdot4\cdot0{,}5\cdot0{,}5=0{,}2$ N. Ötkazgiç induksiya çiziqlariga tik qöyilganda ($\alpha=90^\circ$) kuç eng katta: $F_{max}=BIL=0{,}4$ N.\par\par\medskip\hrule\medskip
\Needspace{12\baselineskip}\noindent\textbf{42.}\ Vakuumda bir-biridan 30 cm masofada $+2\ \mu$C va $+8\ \mu$C nuqtaviy zaryadlar joylaşgan. $k=9\cdot10^9$ N$\cdot$m$^2$/C$^2$.\par
\noindent\hspace*{4mm}a)\ Zaryadlarning özaro ta’sir kuçini (N) toping.\par
\noindent\hspace*{4mm}b)\ Zaryadlarni tutaştiruvçi kesmada maydon kuçlanganligi nolga teng bölgan nuqta birinçi zaryaddan qanday masofada (cm) joylaşgan?\par
\noindent\colorbox{gray!12}{\parbox{\dimexpr\linewidth-2\fboxsep}{\textbf{Javob: a)\ 1,6;\quad b)\ 10}}}\par\smallskip
\noindent\textit{Yechim.}\ \[F=k\frac{q_1q_2}{r^2}=9\cdot10^9\cdot\frac{2\cdot10^{-6}\cdot8\cdot10^{-6}}{0{,}09}=1{,}6\ \text{N}.\]
Bir işorali zaryadlar orasida $kq_1/x^2=kq_2/(r-x)^2$: $(r-x)/x=\sqrt{q_2/q_1}=2$, $x=r/3=10$ cm.\par\par\medskip\hrule\medskip
\Needspace{12\baselineskip}\noindent\textbf{43.}\ Yoruğlik nuri havodan suvga $53^\circ$ tuşiş burçagi ostida tuşmoqda. $n_{suv}=4/3$, $\sin53^\circ=0{,}8$, $c=3\cdot10^8$ m/s.\par
\noindent\hspace*{4mm}a)\ Siniş burçagini (gradusda) toping.\par
\noindent\hspace*{4mm}b)\ Yoruğlikning suvdagi tezligini ($10^8$ m/s) toping.\par
\noindent\colorbox{gray!12}{\parbox{\dimexpr\linewidth-2\fboxsep}{\textbf{Javob: a)\ 37;\quad b)\ 2,25}}}\par\smallskip
\noindent\textit{Yechim.}\ Siniş qonuni: $\sin\beta=\sin53^\circ/n=0{,}8\cdot\frac34=0{,}6$, $\beta=37^\circ$. Suvda $v=c/n=3\cdot10^8\cdot\frac34=2{,}25\cdot10^8$ m/s.\par\par\medskip\hrule\medskip
\Needspace{12\baselineskip}\noindent\textbf{44.}\ Fokus masofasi 10 cm bölgan yiğuvçi linza buyumning 2 marta kattalaşgan haqiqiy tasvirini hosil qildi.\par
\noindent\hspace*{4mm}a)\ Buyumdan linzagaça masofani (cm) toping.\par
\noindent\hspace*{4mm}b)\ Linzadan tasvirgaça masofani (cm) toping.\par
\noindent\colorbox{gray!12}{\parbox{\dimexpr\linewidth-2\fboxsep}{\textbf{Javob: a)\ 15;\quad b)\ 30}}}\par\smallskip
\noindent\textit{Yechim.}\ Haqiqiy tasvir uçun $\Gamma=f/d=2$ va $\frac1d+\frac1f=\frac1F$: $\frac1d+\frac1{2d}=\frac1{10}$, $d=15$ cm, $f=30$ cm.\par\par\medskip\hrule\medskip
\Needspace{12\baselineskip}\noindent\textbf{45.}\ Radioaktiv izotopning yarim yemiriliş davri 8 sutka. Dastlab namunada $1{,}6\cdot10^{20}$ ta radioaktiv yadro bor edi.\par
\noindent\hspace*{4mm}a)\ 24 sutkadan keyin yadrolarning qança uluşi yemirilmay qoladi (önli kasrda)?\par
\noindent\hspace*{4mm}b)\ 24 sutka davomida yemirilgan yadrolar sonini ($10^{20}$) toping.\par
\noindent\colorbox{gray!12}{\parbox{\dimexpr\linewidth-2\fboxsep}{\textbf{Javob: a)\ 0,125;\quad b)\ 1,4}}}\par\smallskip
\noindent\textit{Yechim.}\ 24 sutka — 3 ta yarim yemiriliş davri: $N/N_0=2^{-3}=0{,}125$. Yemirilganlar $\Delta N=1{,}6\cdot10^{20}\cdot0{,}875=1{,}4\cdot10^{20}$.\par\par\medskip\hrule\medskip
\end{document}